% ---------------------------------------------------------------------
%  CHƯƠNG 1 — MỞ ĐẦU
% ---------------------------------------------------------------------
\chapter{Mở đầu}
\label{ch:modau}

\begin{muctieu}
Chương mở đầu trình bày bối cảnh thực tế và lý do chọn đề tài, xác định các
mục tiêu định lượng cần đạt được, giới hạn phạm vi thực hiện và nêu bố cục
tổng thể của báo cáo.
\end{muctieu}

\section{Lý do chọn đề tài}

Trong bối cảnh các hệ thống phân tán và ứng dụng tương tác trực tuyến ngày
càng phổ biến, việc xây dựng một máy chủ mạng có khả năng phục vụ đồng thời
nhiều kết nối khách mà vẫn giữ tính nhất quán tuyệt đối của dữ liệu là kỹ
năng thiết yếu đối với kỹ sư phần mềm. Bài toán \term{Trò chơi đoán số nhiều
người} \en{(Multiplayer Number Guessing Game)} tuy có luật chơi quen thuộc
nhưng hội tụ đầy đủ các thách thức kỹ thuật tiêu biểu của lập trình hệ thống
trên nền tảng Java: giao tiếp mạng hướng dòng qua Socket TCP, điều phối bể
luồng phục vụ nhiều máy khách song song, kiểm soát tranh chấp khi nhiều luồng
cùng đoán trúng ở một thời điểm, luân chuyển lượt chơi bằng hàng đợi, tự động
đóng ván khi hết giờ và quản lý nhiều phòng chơi độc lập.

Đề tài số 30 được nhóm lựa chọn vì tính trực quan, dễ kiểm chứng và bám sát
trọn vẹn sáu chương kiến thức của học phần Lập trình Java nâng cao
(AJP20101)~\cite{dophuchao2026ajp}: từ lập trình hướng đối tượng (Chương 1),
biểu thức Lambda và Streams API (Chương 2), cấu trúc dữ liệu Collections và
so sánh thứ tự (Chương 3), lập trình đa luồng và đồng bộ hóa (Chương 4), vào
ra luồng và lập trình mạng Socket TCP (Chương 5), đến các tính năng hiện đại
của Java 21 như \code{record}, \code{enum}, \code{Optional}, biểu thức
\code{switch} và tính bất biến \en{(immutability)} (Chương 6).

\section{Mục tiêu của đồ án}

Bài tập lớn hướng tới năm mục tiêu cụ thể và đo lường được sau đây:
\begin{itemize}[leftmargin=1.6em, itemsep=2pt]
  \item \textbf{Mục tiêu 1 -- Hạ tầng mạng TCP đa luồng:} Xây dựng máy chủ
        \code{GameServer} theo mô hình \en{thread-per-connection} dùng bể
        luồng cố định \code{Executors.newFixedThreadPool(20)}, phục vụ đồng
        thời tối thiểu 20 máy khách đang hoạt động và xếp hàng an toàn từ máy
        khách thứ 21 trở đi.
  \item \textbf{Mục tiêu 2 -- Giao thức văn bản hướng dòng chuẩn hóa:} Thiết
        kế và cài đặt đầy đủ tập lệnh văn bản theo dòng (\code{JOIN},
        \code{START}, \code{GUESS}, \code{SCORE}, \code{HISTORY},
        \code{STATS}, \code{CREATE\_ROOM}, \code{JOIN\_ROOM},
        \code{LEAVE\_ROOM}, \code{LIST\_ROOMS}, \code{CURRENT\_ROOM},
        \code{QUIT}) cùng các thông điệp phát chung \code{INFO} và gợi ý
        \code{HIGHER}/\code{LOWER}/\code{WIN}.
  \item \textbf{Mục tiêu 3 -- Hoàn thiện 21/21 chức năng (\code{F01}--\code{F21}):}
        Hiện thực đầy đủ cả phần Lõi, phần Phải làm và phần Mở rộng theo bảng
        phân rã chức năng của Đề tài 30, bao gồm cả chế độ đoán đồng thời lẫn
        đoán lần lượt theo vòng (\code{F19}), giới hạn thời gian mỗi ván
        (\code{F20}) và nhiều phòng chơi song song chống bế tắc (\code{F21}).
  \item \textbf{Mục tiêu 4 -- Giao diện máy khách đáp ứng:} Xây dựng cửa sổ
        đồ họa \code{ClientUI} bằng Java Swing tách luồng đọc mạng riêng biệt
        để giao diện không bị treo khi chờ thông điệp từ máy chủ, đồng thời
        hiển thị bảng điểm, lịch sử và thống kê rõ ràng bằng \code{text block}.
  \item \textbf{Mục tiêu 5 -- Kiểm thử tự động TDD đạt 100\%:} Viết kiểm thử
        trước khi viết mã sản xuất cho mọi chức năng, bao phủ toàn bộ trường
        hợp biên và đạt 100\% bài kiểm thử JUnit 5 vượt qua (103/103 tests).
\end{itemize}

\section{Phạm vi và đối tượng}

\textbf{Đối tượng nghiên cứu:} Cơ chế truyền thông mạng hướng dòng TCP
Client--Server, mô hình đồng bộ trạng thái ván chơi đa luồng, thuật toán xoay
vòng lượt đoán, cơ chế hẹn giờ đóng ván và kỹ thuật giành khóa theo thứ tự cố
định giữa nhiều phòng chơi.

\textbf{Phạm vi triển khai:} Hệ thống được phát triển trên nền tảng Java 21
Standard Edition (JDK 21), chỉ sử dụng thư viện chuẩn của JDK và khung kiểm
thử JUnit 5 (\code{junit-platform-console-standalone.jar}). Toàn bộ trạng
thái người chơi trực tuyến, bảng điểm nhiều ván và lịch sử lượt đoán được lưu
trữ trong bộ nhớ bằng các cấu trúc dữ liệu an toàn luồng theo đúng yêu cầu
của đề bài; không sử dụng các thư viện mạng ngoài (như Netty, Grizzly) hay hệ
quản trị cơ sở dữ liệu ngoài phạm vi môn học.

\begin{luuy}
Nguyên tắc kiến trúc xuyên suốt của đồ án là \textbf{tầng nghiệp vụ
(\code{doanso.domain}, \code{doanso.service}) hoàn toàn không phụ thuộc vào
giao diện (\code{doanso.ui}) hay mạng (\code{doanso.net})}. Nhờ đó, toàn bộ
logic ván chơi, đồng bộ 100 luồng và thống kê Streams đều có thể kiểm thử độc
lập với tốc độ mili-giây mà không cần mở cổng mạng.
\end{luuy}

\section{Bố cục báo cáo}

Phần còn lại của báo cáo được tổ chức như sau. Chương~\ref{ch:coso} trình bày
cơ sở lý thuyết và công nghệ sử dụng. Chương~\ref{ch:thietke} phân tích bài
toán, phân rã 21 chức năng và thiết kế kiến trúc 3 tầng. Chương~\ref{ch:caidat}
mô tả quá trình cài đặt các mô-đun cốt lõi, giao diện và bộ kiểm thử TDD.
Chương~\ref{ch:ketqua} trình bày kết quả đạt được, các độ đo đánh giá định
lượng, bàn luận và hạn chế. Chương~\ref{ch:ketluan} tổng kết đồ án, rút ra bài
học kinh nghiệm và nêu hướng phát triển tiếp theo.
